% 导言区
\documentclass[10pt]{article} % 一般只有10 11 12

\usepackage{ctex}

% LaTeX的思想是格式与内容的分离，所以不建议直接在文档中使用大量命令
% 而是用newcommandi命令定义一个新的命令以执行相关的操作
\newcommand{\myfont}{\textif{\textbf{\textsf{}}}}

% 正文区
\begin{document}
    % 字体族设置(罗马字体、无衬线字体、打字机字体)
    % 只设置单独的字体
    \textrm{Roman Family} % 字体族设置命令
    \textsf{Sans Serif Family} % 无衬线字体命令
    \texttt{Typewriter Family} % 打字机字体命令

    % 设置后续的字体
    \rmfamily Roman Family % 字体声明声明后续字体为罗马字体
    % 设置大括号中的字体
    {\sffamily Sans serif Family} % 使用大括号限定作用的范围
    {\ttfamily Typewriter Family}


    % 字体系列设置(粗细、宽度)
    \textmd{Medium Series} \textbf{Boldface Series}

    {\mdseries Medium Series} {\bfseries Boldface Series}

    % 字体形状(直立、斜体、伪斜体、小型大写)
    % 字体设置命令
    \textup{Upright Shape} \textit{Italic Shape}
    \textsl{Slanted Shape} \textsc{Small Caps Shape}

    % 字体设置声明
    {\upshape Upright Shape} {\itshape Italic Shape}
    {\slshape Slanted Shape} {\scshape Small Caps Shape}


    % 中文字体
    {\songti 宋体} \quad {\heiti 黑体} \quad {\fangsong 仿宋} \quad {\kaishu 楷书}

    中文字体的\textbf{粗体}与\textit{斜体}


    % 字体大小
    % 这些声明是以normal set相对大小实现的
    % 而normal set大小是以文档类的参数控制的
    % 文档类的参数是可选参数，可以在方括号中添加
    {\tiny          Hello}\\
    {\scriptsize    Hello}\\
    {\footnotesize  Hello}\\
    {\small         Hello}\\
    {\normalsize    Hello}\\
    {\large         Hello}\\
    {\Large         Hello}\\
    {\LARGE         Hello}\\
    {\huge          Hello}\\
    {\Huge          Hello}\\


    % 中文字号设置命令
    \zihao{-0} 你好！

    \zihao{5} 你好！

    % \myfont

\end{document}
